Define Dedekind domains, give some examples and show that the ring of integers is one.

Definition.
A Dedekind ring (or domain) is a noetherian integrally closed
integral domain \( \mathcal{O} \), where every prime ideal \( \neq (0) \) is a maximal ideal.

Every PID is a Dedekind domain.
Localizations of Dedekind domains are Dedekind
(special case of discrete valuation rings)
Let \( R \) be a Dedekind domain, \( K = \text{Quot}(R)\).
Let \( L \) be a finite degree field extension of \( K \) and denote by \( S \) 
the integral closure of \( R \) in \( L \). Then \( S \) is itself a Dedekind domain.


Proof. 
\(\mathcal{O}_L\) is an integral domain, as it is contained in the field \(L\).
\(\mathcal{O}_L\) is integrally closed by Corollary 2.7(ii).
\(\mathcal{O}_L\) is noetherian: because every ideal \(a \subset \mathcal{O}_L\) is by Corollary 2.21 (0) or generated by 
\(n = [L : \mathbb{Q}]\) elements as a \(\mathbb{Z}\)-module, hence generated by \(\leq n\) elements as an \(\mathcal{O}_L\)-module.

Let \((0) \neq p \subset \mathcal{O}_L\) be a prime ideal \(\implies \mathcal{O}_{L_p}\) is an integral domain and 
\(\#(\mathcal{O}_{L_p}) = (\mathcal{O}_L : p)\) is finite (Theorem 2.23). Every finite integral domain is a field, 
since for every \(0 \neq x \in R\) the map \(R \ni r \mapsto x \cdot r\) is injective \(\underset{R \text{ finite}}{\implies}\) 
bijective \(\implies\) gives us inverse. Hence \(\mathcal{O}_{L_p}\) is a field \(\implies p\) maximal.

The great theorem of Dedekind domains is that their generalize the concept of prime factorizations
in rings which may fail to be UFD.


Prove a lemma that is needed for the prime ideal factorizations of Dedekind domains
and states that (pure) fractional ideals contain the ring of integers.

Lemma 3.8. 
Let \( p \) be a prime ideal in a Dedekind domain \(\mathcal{O}\), let \( L \) be its quotient field \( L = \text{Quot}(\mathcal{O}) \).
(i) For every ideal \( \mathfrak{a} \) there are \( \mathfrak{p_1}, \dots, \mathfrak{p_r} \) with \( \mathfrak{a} \mid \mathfrak{p_1}\cdots\mathfrak{p_r}\)
Let \( p^{-1} := \{ x \in L \mid xp \subset \mathcal{O} \} \).
(i) \( p^{-1} \supseteq \mathcal{O} \).
(ii) For an ideal \( (0) \neq a \subset \mathcal{O} \), let \( ap^{-1} := \{ \sum_{\text{fin}} a_ix_i \mid a_i \in a, x_i \in p^{-1} \} \), then \( ap^{-1} \supseteq a \).

Proof.
(i) Let \( \mathcal{S} := \{ \text{ideals } (0) \neq a \subset \mathcal{O} \text{ without this property} \} \).
Suppose \( \mathcal{S} \neq \emptyset \overset{\mathcal{O} \text{ noeth.}}{\implies} \text{ there is a maximal } a \in \mathcal{S} \).
Observe that \( a \) is not a prime ideal (otherwise \( a \mid p = a \nmid \)), i.e. \(\exists x_1, x_2 \in \mathcal{O} \text{ s.t. } x_1 x_2 \in a, \text{ but } x_1, x_2 \notin a \).
Consider \( a_1 := a + (x_1), a_2 := a + (x_2) \). Both will contain \( a \), but will not be equal to \( a \). By maximality of \( a \), we have \( a_1, a_2 \notin \mathcal{S} \), i.e. \( a_1 \mid p_1 \cdots p_s, a_2 \mid p_{s+1} \cdots p_r \). Then
\(
a \supset (a + (x_1)) \cdot (a + (x_2)) = a_1 \cdot a_2 \supset p_1 \cdots p_s p_{s+1} \cdots p_r \nleq
\)
This is a contradiction, hence \( \mathcal{S} \) must be empty.
(ii) Let \( 0 \neq x \in p \). (i) for \( (x) \) shows \( p \supset (x) \supset p_1 \cdots p_r \), we may choose \( r \) minimal among all such products contained in \( a \).
By Lemma 3.5 \( p \) divides at least one of the factors, e.g. \( p \mid p_i \); 
because prime ideals in Dedekind domains are maximal, \( p_1 = p_1 \).

Now \( (x) \not\supseteq p_2 \cdots p_r \), since \( r \) was minimal, i.e. \(\exists y \in p_2 \cdots p_r\), with \( y \notin (x) \), 
i.e. \(\frac{y}{x} \in L = \text{Quot}(\mathcal{O}) \) does not lie in \(\mathcal{O}\). But since
\( y \cdot p \subset p_2 \cdots p_r \cdot p_r \subset (x) \implies \frac{y}{x} \cdot p \subset \mathcal{O}, \).
so \(\frac{y}{x} \in p^{-1}\).

(iii) Let \( (0) \neq a \subset \mathcal{O} \) be an ideal \(\overset{\mathcal{O} \text{ noeth.}}{\implies}\) \( a \) is finitely generated, e.g. by \(\alpha_1, \ldots, \alpha_n\).
By (ii) we have \( ap^{-1} \supseteq a\mathcal{O} = a \), we need to show "\(\subseteq\)"!
Assume they are equal, then for every \( x \in p^{-1} \) we have \( x\alpha_i \in ap^{-1} = a \) and
\(
x \alpha_i = \sum_{j=1}^{n} a_{ij} x_j \quad \text{with} \quad a_{ij} \in \mathcal{O}.
\)

Then \( A \cdot \begin{pmatrix} \alpha_1 \\ \vdots \\ \alpha_n \end{pmatrix} = 0 \) for \( A = (xI_n - [a_{ij}]) \). Using the Laplace expansion as in Thm. 2.2
\[
\text{det}(A) \cdot \begin{pmatrix} \alpha_1 \\ \vdots \\ \alpha_n \end{pmatrix} = \begin{pmatrix} 0 \\ \vdots \\ 0 \end{pmatrix}
\]
since that is an equation in an integral domain, it must be

\( \text{det}(A) = 0 \).

Thus, \( L \ni x \) is a root of the monic polynomial \(\text{det}(X \cdot I_n - [a_{ij}]) \in \mathcal{O}[X]\), showing that \( x \in \mathcal{O} \), since \(\mathcal{O}\) is integrally closed, \(\forall x \in p^{-1}\).

This would prove that \( p^{-1} \subset \mathcal{O} \) which is absurd as a contradiction to (ii).



Show the prime ideal factorization property of Dedekind domains

Existence:
Let \(\mathcal{S} := \{ a \neq (0), (1) \mid a \text{ does not admit such a decomposition} \}\).
If \( S \neq \emptyset \overset{\mathcal{O} \text{ noeth.}}{\implies} \exists \text{ maximal among the ideals in } \mathcal{S} \). 
Let \( p \) be a maximal ideal containing \( a \) (ex. bc. of Zorn's lemma); 
\( p \supsetneq a \), since \( a \) has no decomposition.

\(
\mathcal{O} \supsetneq p^{-1}pp^{-1} \supsetneq a
\)
and \(p^{-1}pp^{-1} \supsetneq ap^{-1}\), like the last subset inequality, both coming from the Lemma 3.8 (ii).

The vertical (strict!) inclusion implies \( pp^{-1} = \mathcal{O} \), as \( p \) is a maximal ideal.
If \( ap^{-1} = \mathcal{O} \), then
\(
a = a\mathcal{O} = (ap^{-1})p = p.
\)
A contradiction.
Otherwise \(\mathcal{O} \supsetneq ap^{-1} \supsetneq a\), i.e. \( ap^{-1} \notin \mathcal{S} \). 
Consequently, this ideal admits a decomposition \( ap^{-1} = p_1 \cdots p_r \), into prime ideals \(\neq (0) \). But then
\(
a = (ap^{-1})p = p_1 \cdots p_r p.
\)
A contradiction.
In both cases we get a contradiction to \( a \in \mathcal{S} \), thus \( \mathcal{S} = \emptyset \).

Uniqueness: Assume \( a = p_1 \cdots p_r = q_1 \cdots q_s \) \((r \leq s)\).
Then \( p_1 \mid q_1 \cdots q_s \), so by Lemma 3.5 \( p_1 \mid q_i \) 
for at least one \( i \), w.l.o.g. \( i = 1 \). As \( p_1, q_1 \) are maximal ideals, we have \( p_1 = q_1 \).
\(
p_2 \cdots p_r = p_1^{-1} p_1 p_2 \cdots p_r = p_1^{-1} q_2 \cdots q_s = q_2 \cdots q_s.
\)
Repeating this argument \( r \) times (and reordering if needed), we get
\(
p_1 = q_1 \cdots p_r = q_r \quad \text{and} \quad \mathcal{O} = q_{r+1} \cdots q_s \subseteq q_s,
\)
which is a contradiction to \( q_s \) prime, unless \( s = r \). This proves uniqueness.



State the chinese remainder theorem and its Corollary for Dedekind domains

A necessary Lemma.
Let \( a_1, \ldots, a_n \) be ideals in the ring \( R \) with \( a_i + a_j = (1) \) for all \( i \neq j \). Then we have
\(
\bigcap_{i=1}^n a_i = a_1 \cdots a_n.
\)
Proof. 
For each \( i \) we have \( a_i \supset a_1 \cdots a_n \), so \( \bigcap_i a_i \supset a_1 \cdots a_n \).
\( a_i \supset \bigcap_i a_i \Longrightarrow a_1 \cdots a_r \supset \bigcap_i a_i \).
which can be deduced from 3.9 with \(b = \bigcap_i a_i\) and induction

Theorem 3.11 (Chinese remainder theorem). 
Let \( a_1, \ldots, a_n \) be ideals in the ring \( R \) with \( a_i + a_j = (1) \) for all \( i \neq j \). Then
\(
\frac{R}{\bigcap_{i=1}^n a_i} \cong \bigoplus_{i=1}^n \frac{R}{a_i}.
\)

It follows for \(\mathcal{O}\) Dedekind domain, 
\(a = \mathfrak{p}_1^{v_1} \cdots \mathfrak{p}_r^{v_r}\) the prime ideal decomposition of \(a \neq (0)\), that
\(
\frac{\mathcal{O}}{a} \cong \frac{\mathcal{O}}{\mathfrak{p}_1^{v_1}} \oplus \cdots \oplus \frac{\mathcal{O}}{\mathfrak{p}_r^{v_r}},
\)
since \(\mathfrak{p}_i^{v_i} + \mathfrak{p}_j^{v_j} = \mathcal{O}\) (non trivial was an exercise).


Define fractional ideals and show that their form a multiplicative group.

A fractional ideal in \( L \) is a finitely generated \( \mathcal{O} \)-submodule \( \neq 0 \)
of L.

Theorem 3.13. 
The set of fractional ideals in \( L \) together with multiplication forms a group \( J_L \), the ideal group of \( L \). 
The identity is \( (1) = \mathcal{O} \), the inverse of \( \mathfrak{a} \) is 
\(\mathfrak{a}^{-1} := \{ x \in L \mid x \mathfrak{a} \subset \mathcal{O} \} \).

Proof. The operation is clearly associative, commutative and \(\mathfrak{a} \cdot \mathcal{O} = \mathfrak{a}\). 
Furthermore, \(\mathfrak{a}^{-1}\) is a fractional ideal:
\(\mathcal{O}\)-submodule of \( L \): Trivial.
\(\neq (0)\): \(\mathfrak{a}\) is a fractional ideal, hence \(\neq (0)\) and finitely generated. Let \( c \neq 0 \) be the product of denominators 
of some generating set, then \( c \cdot \mathfrak{a} \subset \mathcal{O} \), i.e. \( c \in \mathfrak{a}^{-1} \).
Finitely generated: Let \( 0 \neq y \in \mathfrak{a} \). Then \( y \cdot \mathfrak{a}^{-1} \subset \mathcal{O} \), so by the previous example 
\(\mathfrak{a}^{-1}\) is finitely generated.

We will verify \(\mathfrak{a} \mathfrak{a}^{-1} = \mathcal{O}\) in several steps:
(I) Let \(\mathfrak{p}\) be a prime ideal, as in the proof of 3.6:
\(
\mathcal{O} \supset \mathfrak{p}^{p-1} \Rightarrow \mathfrak{p} \quad \text{and} \quad \mathfrak{p}\mathfrak{p}^{-1} = \mathcal{O}.
\)
(II) For an integral ideal \(\mathfrak{a} = \mathfrak{p}_1 \cdots \mathfrak{p}_r\) (prime ideals \(\neq (0)\)) consider 
\(\mathfrak{b} := \mathfrak{p}_1^{-1} \cdots \mathfrak{p}_r^{-1}\). 
Obviously \(\mathfrak{b}\) is inverse to \(\mathfrak{a}\) because of 
\(\mathfrak{p}_i \mathfrak{p}_i^{-1} = \mathcal{O}\) + associativity + commutativity.
Claim: \(\mathfrak{b} = \mathfrak{a}^{-1}\) (as defined above)
\(\subseteq\): \(\mathfrak{b} \subset \mathfrak{a}^{-1}\) since \(\mathfrak{b} \mathfrak{a} = \mathcal{O}\).
\(\supseteq\): Let \(x \in \mathfrak{a}^{-1}\), by definition \(x \mathfrak{a} \subset \mathcal{O}\). 
Multiplying with \(\mathfrak{b}\) yields \(x (\mathfrak{ab}) \subset \mathfrak{b}\), so \(x \in \mathfrak{b}\). 
This shows \(\mathfrak{a}^{-1} \subset \mathfrak{b}\).
(III) Finally, let \(\mathfrak{a}\) be a fractional ideal, hence \(\exists c \neq 0 : c\mathfrak{a} \subset \mathcal{O}\). 
As seen in (II) previously, the inverse of \(c\mathfrak{a}\) is
\(
(c\mathfrak{a})^{-1} = \{ x \in L \mid x c \mathfrak{a} \subset \mathcal{O} \} = \{ x \in L \mid x c \in \mathfrak{a}^{-1} \} = c^{-1} \mathfrak{a}^{-1}
\).

Using this we can easily verify
\(
\mathfrak{a} \mathfrak{a}^{-1} = (c\mathfrak{a})(c^{-1} \mathfrak{a}^{-1}) = (c\mathfrak{a})(c\mathfrak{a})^{-1} = \mathcal{O}.
\)


State the modified version of the prime ideal decomposition theorem for fractional ideals

Every fractional ideal \( \mathfrak{a} \) in \( L = \text{Quot}(\mathcal{O})\) 
(for a Dedekind domain \( \mathcal{O} \)) can be written uniquely (up to ordering) in the form 
\( \mathfrak{a} = \mathfrak{p}^{v_1}_1 \cdots \mathfrak{p}^{v_r}_r\) for pairwise
distinct prime ideals \( \mathfrak{p}_i \), \( r \geq 0 \) and \( v_i \in \mathbb{Z} \setminus \{0\}\).


Define the ideal class group and fractional principal ideals.

Definition 3.15. 
\( P_L \subset J_L \) is the subgroup of fractional principal ideals in \( L \). 
The quotient \(\text{Cl}_L := J_L / P_L\) is the ideal class group.

It will be later shown that \(\# \text{Cl}_L\) is finite for every number field \( L \).

Remark.
- \( h_L = 1 \Leftrightarrow \# \text{Cl}_L = 1 \Leftrightarrow P_L = J_L \Leftrightarrow \mathcal{O}\) is a PID.
- \( P_L \cong L^\times / \mathcal{O}^\times\), since a generator \(\in L\) of fractional ideal is unique up to units in \(\mathcal{O}^\times\). 
This gives us an exact sequence of abelian groups

\(
1 \longrightarrow \mathcal{O}^\times \longrightarrow L^\times \longrightarrow J_L \longrightarrow \text{Cl}_L \longrightarrow 1
\)


